\chapter{ТЕОРЕТИЧЕСКИЕ ОСНОВЫ DATA-DRIVEN МЕТОДОВ ДЛЯ ДИНАМИЧЕСКИХ СИСТЕМ}
\label{ch:theory}

\section{Постановка задачи восстановления динамики по дискретным данным}
\label{sec:problem_statement}

В большинстве прикладных задач физическое моделирование сложных систем требует значительных вычислительных ресурсов. Альтернативным подходом является использование методов, управляемых данными (data-driven), которые позволяют восстановить оператор эволюции системы исключительно на основе наблюдаемых измерений. 

Пусть нелинейная динамическая система с управляющим воздействием описывается дифференциальным уравнением вида:
\begin{equation}
\label{eq:cont_system}
\frac{d}{dt}x(t) = f(x(t), u(t)),
\end{equation}
где $x(t) \in \mathbb{R}^n$ --- вектор состояния системы, $u(t) \in \mathbb{R}^m$ --- вектор управления. На практике непрерывные траектории недоступны, и система измеряется с заданным шагом дискретизации $\Delta t$. Тогда задача сводится к поиску дискретного оператора сдвига (или его аппроксимации):
\begin{equation}
\label{eq:discrete_system}
x_{k+1} = F(x_k, u_k).
\end{equation}

Для применения матричных алгоритмов данные наблюдений организуются в виде матриц снимков состояний (snapshots). Формируются две матрицы данных $X$ и $X'$, смещенные друг относительно друга на один шаг по времени:
\begin{equation}
\label{eq:matrices_X}
X = \begin{bmatrix} | & | & & | \\ x_1 & x_2 & \dots & x_{N-1} \\ | & | & & | \end{bmatrix}, \quad
X' = \begin{bmatrix} | & | & & | \\ x_2 & x_3 & \dots & x_N \\ | & | & & | \end{bmatrix}.
\end{equation}
Аналогичным образом формируется матрица управляющих воздействий $\Upsilon = [u_1, u_2, \dots, u_{N-1}]$.

\section{Метод Dynamic Mode Decomposition (DMD и DMDc)}
\label{sec:dmd}

\subsection{Классический алгоритм DMD}

Метод динамической декомпозиции мод (Dynamic Mode Decomposition, DMD) аппроксимирует нелинейную динамику локально-линейным оператором $A$. Основная гипотеза DMD заключается в том, что матрицы снимков связаны линейным преобразованием:
\begin{equation}
\label{eq:dmd_base}
X' \approx A X.
\end{equation}
Оптимальная матрица $A$, минимизирующая ошибку аппроксимации в норме Фробениуса, вычисляется через псевдообращение Мура-Пенроуза: $A = X' X^+$. Поскольку размерность вектора состояния $n$ может быть крайне велика, матрица $A$ (размера $n \times n$) вычисляется в редуцированном пространстве. Для этого применяется сингулярное разложение (SVD) матрицы $X \approx U \Sigma V^*$, где $U$ содержит левые сингулярные векторы, $\Sigma$ --- диагональная матрица сингулярных чисел, а $V$ --- правые сингулярные векторы. 

Редуцированная матрица оператора $\tilde{A}$ вычисляется как проекция $A$ на базис $U$:
\begin{equation}
\label{eq:dmd_reduced}
\tilde{A} = U^* A U = U^* X' V \Sigma^{-1}.
\end{equation}
Собственные значения $\tilde{A}$ характеризуют временную эволюцию системы (частоты и затухание мод).

\subsection{DMD с управлением (DMDc)}

Классический DMD не способен корректно извлечь собственную динамику системы, если на нее действует внешнее управление, так как изменения состояния смешиваются с реакцией на входной сигнал. Метод DMDc (DMD with control) расширяет классический подход для систем с управлением \cite{proctor2014dynamic}. Динамика предполагается линейной относительно состояния и управления:
\begin{equation}
\label{eq:dmdc_base}
X' \approx A X + B \Upsilon,
\end{equation}
где $B$ --- матрица управления размером $n \times m$. 

Для нахождения $A$ и $B$ перепишем выражение \ref{eq:dmdc_base} в блочном виде:
\begin{equation}
\label{eq:dmdc_block}
X' \approx \begin{bmatrix} A & B \end{bmatrix} \begin{bmatrix} X \\ \Upsilon \end{bmatrix} = G \Omega.
\end{equation}
К объединенной матрице $\Omega$ применяется сингулярное разложение $\Omega = \tilde{U} \tilde{\Sigma} \tilde{V}^*$. Оператор $G$ аппроксимируется как $G \approx X' \tilde{V} \tilde{\Sigma}^{-1} \tilde{U}^*$. Для разделения матриц собственной динамики $A$ и управления $B$ матрица левых сингулярных векторов $\tilde{U}$ разбивается на блоки, соответствующие размерностям состояний и управления, после чего применяются методы проекции на базис наблюдений без учета влияния входов \cite{proctor2014dynamic}.


\section{Метод Sparse Identification of Nonlinear Dynamics (SINDy и SINDy-c)}
\label{sec:sindy}

В отличие от DMD, который ищет линейный оператор эволюции, метод SINDy направлен на поиск явного нелинейного дифференциального (или разностного) уравнения, управляющего системой \cite{brunton2016discovering}. Основная идея SINDy заключается в том, что большинство физических систем описывается небольшим числом активных слагаемых из обширного пространства возможных нелинейных функций.

Для непрерывной системы $\dot{x} = f(x, u)$ формируется библиотека кандидатных функций $\Theta(X, \Upsilon)$. Матрица $\Theta$ состоит из столбцов, представляющих собой нелинейные комбинации состояний и управлений (например, полиномы, тригонометрические функции):
\begin{equation}
\label{eq:sindy_library}
\Theta(X, \Upsilon) = \begin{bmatrix} 1 & X & \Upsilon & X \otimes X & X \otimes \Upsilon & \sin(X) & \dots \end{bmatrix}.
\end{equation}

Тогда производные по времени (вычисленные численно по данным $X$) можно представить в виде линейной комбинации функций из библиотеки:
\begin{equation}
\label{eq:sindy_approx}
\dot{X} \approx \Theta(X, \Upsilon) \Xi,
\end{equation}
где $\Xi = [\xi_1, \xi_2, \dots, \xi_n]$ --- матрица коэффициентов. Каждый столбец $\xi_i$ определяет, какие нелинейные функции формируют динамику $i$-й переменной состояния.

Поскольку искомое уравнение предполагается разреженным (sparse), задача сводится к поиску матрицы $\Xi$ с минимальным числом ненулевых элементов. Это достигается путем решения оптимизационной задачи с $L_1$-регуляризацией или с использованием алгоритма STLSQ (Sequential Thresholded Least Squares):
\begin{equation}
\label{eq:stlsq}
\xi_i = \arg \min_{\xi_i} \| \dot{X}_i - \Theta(X, \Upsilon)\xi_i \|_2 + \lambda \| \xi_i \|_0,
\end{equation}
где $\lambda$ --- порог отсечения. Элементы матрицы $\Xi$, абсолютное значение которых меньше $\lambda$, принудительно обнуляются на каждой итерации, а оставшиеся коэффициенты пересчитываются методом наименьших квадратов.

Наличие управления $\Upsilon$ (SINDy-c) естественным образом интегрируется в метод путем добавления управляющих переменных в словари $\Theta$, что позволяет алгоритму находить перекрестные нелинейные члены вида $x_i u_j$.